% Einheit 2 – Normalform und p-q-Formel – Hauptnummern 24 bis 39
\einheitenkopf[e2][2 p-q-Formel]{Einheit 2 von 4 · Normalform und p-q-Formel}
\zweigzeile{Hier lernst du, jede quadratische Gleichung zu ordnen und mit der p-q-Formel zu lösen · neu in diesem Jahr, je nach Buch bis Klasse 10 · P10 oft · baut auf: Wurzelziehen (Einheit 1), Ausmultiplizieren, Terme ordnen}
\setcounter{aufgabe}{23}

\begin{aufgabe}{Ich erkenne an der Form einer Gleichung, welcher Lösungsweg passt.}
Kreuze den passenden Lösungsweg an. Löse nichts.
\begin{teile}
\teil $x^2 - 7x + 10 = 0$ \hfill \kreuz{Wurzelziehen} \kreuz{rückwärts rechnen} \kreuz{p-q-Formel}
\teil $x^2 = 30$ \hfill \kreuz{Wurzelziehen} \kreuz{rückwärts rechnen} \kreuz{p-q-Formel}
\teil $(x - 4)^2 = 5$ \hfill \kreuz{Wurzelziehen} \kreuz{rückwärts rechnen} \kreuz{p-q-Formel}
\teil $3x^2 - 12 = 0$ \hfill \kreuz{Wurzelziehen} \kreuz{rückwärts rechnen} \kreuz{p-q-Formel}
\teil $x^2 + x = 6$ \hfill \kreuz{Wurzelziehen} \kreuz{rückwärts rechnen} \kreuz{p-q-Formel}
\end{teile}
\end{aufgabe}

\begin{aufgabe}{Ich erkenne p und q in einer Gleichung in Normalform.}
Kreise p ein und unterstreiche q, jeweils mit dem Vorzeichen davor. Rechne nichts.
\begin{teilezwei}
\tz $x^2 + 7x + 2 = 0$ & \tz $x^2 - 2x + 9 = 0$ \\
\tz $x^2 + 5x - 1 = 0$ & \tz $x^2 - x - 20 = 0$ \\
\tz $x^2 - 11 = 0$ & \\
\end{teilezwei}
\end{aufgabe}

\begin{aufgabe}{Ich erkenne, ob vor $x^2$ eine Eins steht.}
Kreuze an, ob du vor der p-q-Formel teilen musst. Unterstreiche die Zahl, durch die du teilen musst.
\begin{teile}
\teil $x^2 + 5x - 3 = 0$ \hfill \kreuz{teilen} \kreuz{nicht teilen}
\teil $2x^2 + 4x - 7 = 0$ \hfill \kreuz{teilen} \kreuz{nicht teilen}
\teil $-x^2 + 3x + 4 = 0$ \hfill \kreuz{teilen} \kreuz{nicht teilen}
\teil $0{,}5x^2 - x + 2 = 0$ \hfill \kreuz{teilen} \kreuz{nicht teilen}
\teil $x^2 - 9x + 1 = 0$ \hfill \kreuz{teilen} \kreuz{nicht teilen}
\end{teile}
\end{aufgabe}

\verfahren{Lösen mit der p-q-Formel}
\begin{aufgabe}{Ich kann eine Gleichung in Normalform mit der p-q-Formel lösen. (je nach Buch erst in Klasse 10)}
Löse mit der p-q-Formel.
\rechenplatz{4}
\begin{gleichungsraster}[3]
\gl{x^2 + 6x + 8 = 0} & \gl{x^2 + 10x + 16 = 0} \\
\gl{x^2 + 6x + 5 = 0} & \gl{x^2 + 10x + 21 = 0} \\
\end{gleichungsraster}
\end{aufgabe}

\begin{aufgabe}{Ich kann eine Gleichung in Normalform mit der p-q-Formel lösen – weiter.}
\begin{gleichungsraster}[3]
\gl{x^2 - 4x - 12 = 0} & \gl{x^2 + 3x - 10 = 0} \\
\anweisung{Gib die Lösungen mit Wurzel und als Näherungswerte auf zwei Stellen nach dem Komma an.}
\gl{x^2 - 4x + 1 = 0} & \gl{x^2 + x - 3 = 0} \\
\anweisung{Berechne die Nullstellen der Parabel. (P10 2025 OS)}
\gl{y = x^2 - 8x + 14} & \\
\end{gleichungsraster}
\end{aufgabe}

\verfahren{Gleichungen ordnen und durch die Zahl vor $x^2$ teilen}
\begin{aufgabe}{Ich kann eine Gleichung erst ordnen und dann mit der p-q-Formel lösen.}
Löse die Gleichung.
\rechenplatz{5}
\begin{gleichungsraster}[3]
\gl{x^2 + 2x = 8} & \gl{x^2 = 6x - 5} \\
\gl{x^2 + 4 = 5x} & \gl{x^2 + 9x - 5 = 3x + 11} \\
\end{gleichungsraster}
\end{aufgabe}

\begin{aufgabe}{Ich kann eine Gleichung mit einer Zahl vor $x^2$ lösen.}
Löse die Gleichung.
\rechenplatz{5}
\begin{gleichungsraster}[3]
\gl{2x^2 - 4x - 6 = 0} & \gl{3x^2 + 9x - 12 = 0} \\
\gl{-x^2 + 6x - 8 = 0} & \gl{0{,}5x^2 - x - 4 = 0} \\
\anweisung{Berechne die Nullstellen der Parabel. (P10 2020 OS)}
\gl{y = 3x^2 - 3x - 36} & \\
\end{gleichungsraster}
\end{aufgabe}

\begin{aufgabe}{Ich kann eine Gleichung mit Klammern erst auflösen, ordnen und dann lösen.}
Löse die Gleichung.
\rechenplatz{5}
\begin{gleichungsraster}[3]
\gl{x \cdot (x + 4) = 12} & \gl{(x + 1)^2 = 3x + 7} \\
\gl{(x - 3)^2 - 1 = -2x + 8} & \\
\anweisung{Für welche x hat die Parabel den y-Wert $-8$? (P10 2023 OS)}
\gl{y = -x^2 + 6x - 1} & \\
\end{gleichungsraster}
\end{aufgabe}

\verfahren{Zahl der Lösungen}
\begin{aufgabe}{Ich kann eine Gleichung mit genau einer Lösung lösen.}
Löse mit der p-q-Formel.
\begin{gleichungsraster}[3]
\gl{x^2 + 6x + 9 = 0} & \gl{x^2 - 10x + 25 = 0} \\
\gl{x^2 - x + 0{,}25 = 0} & \gl{4x^2 + 12x + 9 = 0} \\
\end{gleichungsraster}
\end{aufgabe}

\begin{aufgabe}{Ich kann erkennen, dass eine Gleichung keine Lösung hat.}
Löse mit der p-q-Formel. Hat die Gleichung keine Lösung, schreibe das als Satz.
\begin{gleichungsraster}[3]
\gl{x^2 + 2x + 5 = 0} & \gl{x^2 - 6x + 10 = 0} \\
\gl{x^2 + x + 1 = 0} & \gl{2x^2 + 4x + 7 = 0} \\
\end{gleichungsraster}
\end{aufgabe}

\begin{aufgabe}{Ich kann am Wert unter der Wurzel entscheiden, wie viele Lösungen es gibt.}
Berechne nur den Wert unter der Wurzel und kreuze an.
\begin{teile}
\teil $x^2 + 4x + 1 = 0$ \hfill Wert: \leerfeld \quad \kreuz{zwei} \kreuz{eine} \kreuz{keine}
\teil $x^2 - 12x + 36 = 0$ \hfill Wert: \leerfeld \quad \kreuz{zwei} \kreuz{eine} \kreuz{keine}
\teil $x^2 + 2x + 4 = 0$ \hfill Wert: \leerfeld \quad \kreuz{zwei} \kreuz{eine} \kreuz{keine}
\teil $x^2 - 3x + 1 = 0$ \hfill Wert: \leerfeld \quad \kreuz{zwei} \kreuz{eine} \kreuz{keine}
\teil $x^2 + 5x + 7 = 0$ \hfill Wert: \leerfeld \quad \kreuz{zwei} \kreuz{eine} \kreuz{keine}
\anweisung{Den Wert unter der Wurzel nennt man auch Diskriminante.}
\teil $x^2 - 2x - 1 = 0$ \hfill Diskriminante: \leerfeld \quad \kreuz{zwei} \kreuz{eine} \kreuz{keine}
\end{teile}
\end{aufgabe}

\verfahren{Probe und Lösungsweg}
\begin{aufgabe}{Ich kann mit einer Probe prüfen, ob eine Lösung stimmt.}
Setze die Zahl ein und kreuze an: wahre Aussage (wA) oder falsche Aussage (fA).
\begin{teile}
\teil $x = 2$ \quad in \quad $x^2 + 2x - 8 = 0$ \hfill \kreuz{wA} \kreuz{fA}
\teil $x = -1$ \quad in \quad $x^2 - 5x + 4 = 0$ \hfill \kreuz{wA} \kreuz{fA}
\teil $x = -1$ \quad in \quad $x^2 - 2x - 3 = 0$ \hfill \kreuz{wA} \kreuz{fA}
\teil $x = 3$ \quad in \quad $x^2 + x - 6 = 0$ \hfill \kreuz{wA} \kreuz{fA}
\teil $x = -4$ \quad in \quad $x^2 + 3x - 4 = 0$ \hfill \kreuz{wA} \kreuz{fA}
\end{teile}
\end{aufgabe}

\begin{aufgabe}{Ich kann den kürzesten Lösungsweg wählen.}
Löse die Gleichung auf dem kürzesten Weg.
\begin{gleichungsraster}[3]
\gl{x^2 - 20 = 5} & \gl{x^2 + 2x - 24 = 0} \\
\gl{(x + 3)^2 = 16} & \gl{x^2 - 6x = 16} \\
\end{gleichungsraster}
\end{aufgabe}

\begin{aufgabe}{Ich finde den Fehler beim Lösen mit der p-q-Formel.}
Finde den Fehler, benenne ihn und korrigiere die Rechnung.
\begin{teile}
\teil Mia rechnet: \rechnung{2x^2 + 16x + 14 &= 0 \qquad p = 16,\ q = 14 \\ x_{1,2} &= -8 \pm \sqrt{64 - 14} = -8 \pm \sqrt{50}}
\teil Tim rechnet: \rechnung{x^2 - 10x + 9 &= 0 \\ x_{1,2} &= -5 \pm \sqrt{25 - 9} = -5 \pm 4 \\ x_1 &= -1, \quad x_2 = -9}
\teil Ella rechnet: \rechnung{x^2 - 8x + 12 &= 0 \\ x_{1,2} &= 4 \pm \sqrt{16 + 12} = 4 \pm \sqrt{28}}
\end{teile}
\end{aufgabe}

\begin{aufgabe}{Ich kann begründen, wann die p-q-Formel funktioniert.}
Begründe.
\begin{teile}
\teil Warum darf man die p-q-Formel bei $2x^2 + 6x - 8 = 0$ nicht sofort anwenden?
\teil Warum hat eine Gleichung keine Lösung, wenn der Wert unter der Wurzel negativ ist?
\end{teile}
\end{aufgabe}

\begin{aufgabe}{Ich kann berechnen, wo sich eine Gerade und eine Parabel schneiden.}
Gezeichnet sind die Parabel zu $y = x^2 - 2x - 1$ und die Gerade zu $y = x + 3$.

\begin{ksys}[xmin=-3,xmax=6,ymin=-3,ymax=8,ablesen]
\parabel{1}{1}{-2}{}
\gerade[5]{1}{3}{}
\end{ksys}

\begin{teile}
\teil Setze gleich und berechne die x-Werte der Schnittpunkte.
\teil Berechne die y-Werte und gib die Schnittpunkte an.
\teil Kontrolliere deine Schnittpunkte am Graphen. \hfill (P10 2024 OS)
\end{teile}
\end{aufgabe}
